\chapter{Kinetic and stochastic simulation}
Chapter 27 determined which molecular-count states a symbolic reaction network
could reach. A reaction edge did not specify a waiting time. This chapter adds
time without confusing three objects: a concentration trajectory, one random
\Term{count trajectory}, and a probability distribution over count trajectories.

Our reference is deliberately small: the reversible conversion $A\rightleftharpoons B$.
Its first-order kinetics admit independent closed-form mean and variance oracles.
That makes it a useful laboratory for numerical reasoning before a larger DNA
circuit obscures errors. The symbols denote abstract species, not validated strand
sequences. No kinetic constant below is an experimentally fitted DNA rate.

\section{One network, two timed semantics}
Write the two reaction channels and their count changes as
\[
R_f:A\longrightarrow B,\quad \nu_f=(-1,1)^\top,\qquad
R_r:B\longrightarrow A,\quad \nu_r=(1,-1)^\top.
\]
Every event preserves $A+B=N$. A rate constant adds a rule for how frequently an
enabled channel occurs; it does not change its stoichiometry. First-order constants
$k_f,k_r$ have units inverse time.

\Plate{01-semantics}{Stoichiometry is shared; state representation and timing law
are not. Continuous concentrations and discrete molecule counts require different
interpretations of the same reaction channels.}{dnad28:semantics}
For concentrations $a,b$ in mol/L, mass action gives
\[
\dot a=-k_fa+k_rb,\qquad \dot b=k_fa-k_rb.
\]
For integer counts $A,B$, the stochastic \Term{propensities} are
$\alpha_f=k_f A,\alpha_r=k_r B$, in events per time. In a well-mixed volume $V$,
$\Omega=N_{\rm A}V$ converts concentration to count: $A=\Omega a$. For this
first-order network, multiplying the concentration equation by $\Omega$ yields
the same algebra for continuous copy-equivalents.

A continuous value such as 26.57 copy-equivalents is not a fractional physical
molecule. It can represent an ensemble mean or a scaled concentration. A single
count path stays integer-valued. The equivalence of their means here follows from
linear propensities, not from an unrestricted rule that ODEs reproduce all
stochastic chemistry.

\section{Derive the deterministic oracle first}
Conservation implies $b=C-a$, where $C=a_0+b_0$. Let
$\lambda=k_f+k_r$. For $\lambda>0$,
\[
\dot a=-\lambda a+k_rC,\qquad
a_*=\frac{k_rC}{\lambda},\qquad
a(t)=a_*+(a_0-a_*)e^{-\lambda t}.
\]
The conserved mode has eigenvalue zero; the deviation from equilibrium has
eigenvalue $-\lambda$. Thus equilibrium and relaxation time are separate facts:
$k_r/\lambda$ fixes the equilibrium fraction in $A$, while $1/\lambda$ fixes the
relaxation scale.

If both rates vanish, the state stays at its initial value. Dividing by
$\lambda$ without this branch produces an invalid calculation at a legitimate
physical boundary. For very small $t$, directly subtracting
$1-e^{-\lambda t}$ loses relative precision. The reference instead evaluates
$q=-\operatorname{expm1}(-\lambda t)$ and writes
\[
a(t)=a_0(1-k_fq/\lambda)+b_0 k_rq/\lambda.
\]
Tests compare this transition-probability expression with the independently
derived equilibrium expression, including zero rates and tiny times.

\section{Euler's method: conservation, positivity, and accuracy}
Explicit Euler updates both species from the same old state:
\[
\begin{pmatrix}a_{n+1}\\b_{n+1}\end{pmatrix}
=
\begin{pmatrix}
1-hk_f&hk_r\\ hk_f&1-hk_r
\end{pmatrix}
\begin{pmatrix}a_n\\b_n\end{pmatrix}.
\]
The columns sum to one, so total population is preserved in exact arithmetic.
Sequentially updating $a$ and then computing the new $b$ flux from the new $a$
would implement a different method.

\Listing{euler}
The transition matrix is nonnegative for every nonnegative initial state exactly
when $h\max(k_f,k_r)\le1$. This is a \Term{positivity bound}. The decaying
mode has multiplier $1-h\lambda$, so linear absolute stability requires
$|1-h\lambda|\le1$, or $h\lambda\le2$. A method can satisfy stability yet
produce negative populations: with $(k_f,k_r)=(0.4,0.1)$ and $h=3$,
the multiplier is $-0.5$, but an initial $(a,b)=(1,0)$ gives $a_1=-0.2$.
The reference rejects this step rather than clipping a negative value and hiding
the violated contract. Equality at the stability boundary permits an undamped
oscillation; stability is not convergence to equilibrium.

\Plate{03-ode}{The deterministic oracle is a smooth ensemble mean. The numerical
table checks first-order convergence; the positivity and stability bounds answer
different questions.}{dnad28:ode}
For $(A_0,B_0)=(100,0)$, rates $(0.4,0.1)$ and $T=5$, the exact mean is
$26.566799889911906$. Euler's errors are generated from the tested reference:
\begin{center}\input{\Generated/euler.tex}\end{center}
Halving the step approximately halves the terminal error. A stronger independent
oracle is $a_n=20+80(1-0.5h)^n$, valid for this fixture. It checks every time
step, not only a terminal comparison against a shared helper.

\section{Derive the direct stochastic simulation method}
At a fixed state $x$, set $\alpha_0=\alpha_f+\alpha_r$. Under the continuous-time
Markov chain model, the probability that no event occurs before time $\tau$ is
$S(\tau)=e^{-\alpha_0\tau}$. Inverse-transform sampling gives
\[
\tau=-\log U_1/\alpha_0,\qquad
\Pr(J=j\mid x)=\alpha_j/\alpha_0.
\]
The joint density is $\alpha_j e^{-\alpha_0\tau}$; the waiting time and channel
choice are conditionally independent given the state. This is the direct-method
construction described in the selected primary exposition
\cite{higham2008}; Gillespie's original paper establishes the historical
stochastic-simulation formulation \cite{gillespie1977}.

\Listing{event_draw}
The first uniform lies strictly in $(0,1)$; the channel uniform lies in $[0,1)$.
The cumulative threshold assigns an exact boundary to the channel on its right.
Clamping a zero waiting-time uniform to an arbitrary epsilon changes the tail of
the waiting-time law. The reference redraws the excluded endpoint instead.
If $\alpha_0=0$, the state is absorbing and no logarithm or division is needed.

\Plate{02-ssa}{Two random draws have distinct roles. Update the count state only
after a completed event; a path stopped by its event budget must not masquerade
as an observation at the requested horizon.}{dnad28:ssa}
\Listing{ssa}
The result records events, requested horizon, simulated interval, completion flag,
and stopping reason. A sampled event after the horizon need not be applied: the
state before it remains valid through the horizon. A zero-rate absorbing state
also completes the interval. Exhausting the computational event budget instead
returns a \Term{censored trajectory}, not a completed horizon observation.

The stored clock must advance strictly and remain finite at accepted events.
When floating-point resolution cannot represent a later event time, failure is
preferable to silently entering an infinite loop. These guards do not make the
simulation exact arithmetic: exact SSA refers to sampling the declared CTMC
without a fixed time-discretization approximation, subject to arithmetic and
random-number implementation.

\section{A path is a right-continuous staircase}
At an event time, the observed count is the post-event state. Between event times,
it stays constant. A straight line connecting event values incorrectly suggests
fractional changes between events. The authored path figure uses a step plot.

The fixed seed-28 path has 111 events by $T=5$ and ends at $(29,71)$.
Its final $A$ count differs from the mean by about 2.43; that is not evidence of a
bug. One sample is supposed to fluctuate. Every transition changes counts by
$(-1,1)$ or $(1,-1)$, and every stored state has nonnegative integer counts totaling
100. The seed is local: importing or running this reference does not reset the
caller's global random-number generator.

A query after an incomplete path's simulated interval is rejected. Reusing its
last count as if it were the horizon count can bias ensemble estimates, because
high-activity trajectories are more likely to hit an event cap. Discarding such
paths can introduce another selection bias. Report interruption frequency and
resolve the budget or use a justified alternative method before estimating a
horizon distribution.

\section{An exact finite-copy distribution}
Each molecule independently switches between two states with the declared
first-order rates. With $q=1-e^{-\lambda t}$,
\[
p_{AA}=1-\frac{k_f}{\lambda}q,\qquad
p_{BA}=\frac{k_r}{\lambda}q.
\]
Initially-$A$ molecules contribute a binomial group with parameter $p_{AA}$;
initially-$B$ molecules contribute another independent group with parameter $p_{BA}$.
Therefore
\[
\mathbb E[A(t)]=A_0p_{AA}+B_0p_{BA},
\quad
\operatorname{Var}A(t)=A_0p_{AA}(1-p_{AA})+B_0p_{BA}(1-p_{BA}).
\]
\Listing{exact_moments}
For a mixed initial population, the sum is not generally a single binomial with
success probability $\mathbb E[A]/N$. The tests enumerate all binary configurations
of a five-molecule mixed fixture and compute its law, mean, and variance without
calling SSA or the moment helper.

\Plate{04-ensemble}{One stochastic path, an exact population distribution, and a
Monte Carlo mean answer different questions. The error bar is uncertainty in the
estimated mean, not the biological spread of individual count outcomes.}{dnad28:ensemble}
For the all-$A$ initial fixture, $A(T)\sim\operatorname{Binomial}(100,p_{AA})$,
with variance $19.508851326005672$ and population standard deviation about 4.417.
The 2,000 fixed-seed trajectories give mean 26.702 and sample variance
19.51895547773887. The estimated standard error of the mean is about 0.09879.
Its discrepancy from the exact mean is about 1.37 exact-law standard errors.

The sample variance estimates the spread of individual paths. Dividing it by the
sample count estimates the variance of the sample mean under independent sampling.
More trajectories shrink uncertainty about the mean; they do not shrink the
physical count variance of this model. Deterministic pseudorandom seeds support
reproducibility, not a mathematical proof of random-number independence or
experimental validation.

\section{When first moments do not close}
For any count-network generator with stoichiometric columns $\nu_j$,
\[
\frac{d}{dt}\mathbb E[X]=\sum_j\nu_j\mathbb E[\alpha_j(X)].
\]
Linear propensities permit replacing the expectation of each propensity with its
value at the mean. Nonlinear propensities generally do not. If $N$ is zero or two
with equal probability, $\mathbb E[N]=1$ but
$\mathbb E[N(N-1)]=1$, whereas $\mathbb E[N](\mathbb E[N]-1)=0$.
This counterexample is enough to invalidate an unqualified mean-field substitution.
Moment closures introduce additional assumptions; increasing molecule count alone
does not prove them for every reaction network or time horizon.

A larger DNA circuit may also invalidate the well-mixed CTMC approximation through
spatial encounters, surface attachment, temperature changes, hidden complexes,
or time-varying reagents. A model should declare whether those effects are omitted,
lumped into effective constants, or resolved explicitly.

\section{Bimolecular units and the identical-reactant factor}
For distinct species with molar event-rate law $v=k\,a\,b$,
\[
\alpha_{AB}(A,B)=\frac{k}{\Omega}AB,\qquad
\Omega=N_{\rm A}V.
\]
Here $k$ has units L/(mol$\cdot$time); multiplying $v$ by $\Omega$ converts
molar concentration change rate into events per time.

For $2A\to P$, declare the convention before writing the factor.
If the molar \emph{event-rate} law is $v=k a^2$, then
\[
\alpha_{AA}(A)=\frac{k}{\Omega}A(A-1)
=\frac{2k}{\Omega}\binom A2.
\]
Each event consumes two $A$ molecules, so $\dot a=-2v$.
If instead a reported constant defines $\dot a=-\widetilde k a^2$, then
$k=\widetilde k/2$. Mixing these conventions creates a factor-of-two error.
The finite-copy propensity approaches the declared mass-action law at large
copy count; its exact pair factor also disables the reaction at count one
\cite{higham2008}.

\Plate{05-units}{Count pairs, concentration scaling, and event stoichiometry are
three separate operations. The identical-reactant convention is stated explicitly
rather than hidden inside a combinatorial factor.}{dnad28:units}
\Listing{bimolecular_propensity}
For $k=10^6$ L/(mol$\cdot$s), $A=10$, $B=20$, $V=10^{-15}$ L,
the distinct-reactant propensity is about 0.332108 events/s. Doubling the volume
at fixed molecule counts halves it. Doubling volume while keeping concentrations
fixed doubles counts and approximately doubles the event rate. These are
different physical interventions, so volume dependence must state what is held fixed.

\section{Validation before calibration}
\Plate{06-validation}{Validation gates connect exact laws to numerical evidence
without upgrading a model check into an experimental chemistry claim.}{dnad28:validation}
The reference checks exact count types, finite nonnegative rates and times,
positive volume and Euler step size, finite propensities, conservation, enabled
unit jumps, stopping status, and independent analytic laws. Booleans and fractional
counts are rejected instead of being silently converted to integers.
Finite inputs alone do not guarantee all intermediate operations remain finite;
the code explicitly checks key outputs but does not promise arbitrary-magnitude
numerical robustness.

This chapter's 35 tests include counterexamples rather than only agreement tests.
No wall-clock benchmark, fitted rate, sequence-level thermodynamic model, or wet-lab
readout is claimed. Chapter 29's differentiable molecular and sequence models will
differentiate declared kinetics and discuss sensitivities. Passing today's
simulation tests is a prerequisite for that work, not evidence that the chosen
rates describe a real assay.

\section{Worked exercises}
\begin{enumerate}
\setlength{\itemsep}{1pt}
\item Derive the conserved quantity. \emph{Solution:} adding the two ODEs gives
$\dot a+\dot b=0$; either count event also changes $A+B$ by zero.
\item Find equilibrium for the fixture. \emph{Solution:} $a_*=0.1(100)/0.5=20$.
\item Find the relaxation time. \emph{Solution:} $1/(0.4+0.1)=2$ time units.
\item Derive Euler's scalar oracle. \emph{Solution:} subtract equilibrium to obtain
$a_{n+1}-20=(1-0.5h)(a_n-20)$ and iterate.
\item Find the largest universally positive step. \emph{Solution:}
$h\le1/\max(0.4,0.1)=2.5$.
\item Find the stability endpoint. \emph{Solution:} $h\le2/0.5=4$;
at four the deviation alternates without damping.
\item Exhibit a stable negative step. \emph{Solution:} at $h=3$ the decaying
multiplier is $-0.5$, yet $(1,0)$ maps to $(-0.2,1.2)$.
\item Why not clip this result? \emph{Solution:} clipping changes the update and
can change conservation unless another ad hoc adjustment is introduced.
\item Draw a waiting time for total propensity four and $U_1=e^{-2}$.
\emph{Solution:} $\tau=2/4=0.5$.
\item Choose a channel with $(\alpha_f,\alpha_r)=(3,1)$ and $U_2=0.75$.
\emph{Solution:} the right-boundary convention selects the reverse channel.
\item Interpret zero total propensity. \emph{Solution:} no enabled timed event;
the state is absorbing over the remaining interval under constant parameters.
\item Query a path at an event time. \emph{Solution:} take its post-event count,
not the previous count or a linearly interpolated value.
\item What does an event-budget stop prove? \emph{Solution:} only the recorded
prefix; it does not supply the requested horizon observation.
\item Derive the all-$A$ variance. \emph{Solution:} sum 100 independent Bernoulli
variances to obtain $100p_{AA}(1-p_{AA})$.
\item Why is the mixed initial law different? \emph{Solution:} the two independent
groups have different transition probabilities before equilibrium.
\item Compare standard deviation with mean standard error. \emph{Solution:}
the former describes one-path spread; the latter divides that spread by $\sqrt M$.
\item Falsify nonlinear mean closure. \emph{Solution:} the zero-or-two example
has factorial second moment one but plug-in factorial moment zero.
\item Convert an identical-reactant event constant. \emph{Solution:}
the coefficient multiplying $\binom A2$ is $2k/\Omega$ for $v=ka^2$.
\item Hold copies fixed and double volume. \emph{Solution:} $\Omega$ doubles,
so a bimolecular propensity halves.
\item Specify a stronger validation experiment. \emph{Solution:} compare fitted
and held-out time-course distributions with declared measurement noise and
replicate-level uncertainty; do not reuse calibration data as an independent test.
\end{enumerate}
